\documentclass[11pt]{article}

% ---- packages ----
\usepackage[utf8]{inputenc}
\usepackage[T1]{fontenc}
\usepackage{amsmath,amssymb}
\usepackage{graphicx}
\usepackage{booktabs}
\usepackage{hyperref}
\usepackage{xcolor}
\usepackage{subcaption}
\usepackage[margin=1in]{geometry}
\usepackage{natbib}

% ---- helper for plan notes ----
\newcommand{\todo}[1]{\textcolor{red}{\textbf{[TODO: #1]}}}
\newcommand{\plan}[1]{\textcolor{blue}{\textit{#1}}}

\title{Yedi: A Multi-Dimensional Strategic Reasoning Benchmark\\for Large Language and Vision-Language Models}

\author{
  \todo{Authors and affiliations}
}

\date{}

\begin{document}
\maketitle

% ======================================================================
\begin{abstract}
\plan{%
  3--4 sentences. Key points to hit:
  \begin{itemize}
    \item Existing game benchmarks test single-axis reasoning (math, spatial, language).
      Real decisions require simultaneous optimization across multiple cognitive
      dimensions under resource constraints.
    \item We present Yedi, a configurable card-game benchmark where merge outcomes
      are determined by multiplicative interaction across 1--4 independent
      dimensions (numerical, visual, spatial, verbal). One bad dimension
      ruins the compound gain --- models must optimize the worst-case factor,
      not the average.
    \item We provide a Gymnasium environment with 40+ difficulty configurations,
      metadata and screenshot observation modes, and three baseline agents
      (random, greedy, LLM/VLM).
    \item Early results: frontier LLMs surpass a strong greedy baseline in
      metadata mode. Comparing metadata vs.\ vision modes on identical
      configs quantifies a ``vision tax'' that isolates perceptual grounding
      cost from strategic reasoning ability.
  \end{itemize}
}
\end{abstract}

% ======================================================================
\section{Introduction}
\label{sec:intro}

\plan{%
  Motivation paragraph (why another benchmark?):
  \begin{itemize}
    \item LLMs/VLMs are increasingly evaluated on game-playing tasks
      (Tetris, 2048, chess, Go, Atari). These are valuable but typically
      test a single axis --- spatial planning, arithmetic, pattern recognition.
    \item Real-world decisions (medical, financial, engineering) require
      balancing multiple incommensurable factors simultaneously. A treatment
      that is optimal for one organ system but harmful to another is not a
      good treatment. Current benchmarks do not test this compound trade-off
      reasoning.
    \item Games with multi-dimensional scoring are a natural testbed, but
      most such games are too complex to evaluate systematically
      (Civilization, Factorio) or too simple to differentiate models
      (multi-objective toy gridworlds).
  \end{itemize}
}

\plan{%
  What Yedi offers (contribution paragraph):
  \begin{itemize}
    \item A card game with 4 independent scoring dimensions (number, color,
      shape, word), each with multiple sub-modes. Merge outcomes are the
      PRODUCT of per-dimension multipliers, so one bad factor dominates.
    \item Dual resource constraints (mana budget + move budget) that force
      planning beyond greedy single-step optimization.
    \item A configurable difficulty surface: 40+ mode combinations from
      single-dimension (isolate one cognitive skill) to 4-dimension
      (compound trade-off reasoning).
    \item Two observation modes --- structured metadata and raw screenshots ---
      on the same underlying task, enabling controlled measurement of the
      perceptual grounding cost (``vision tax'').
    \item Open-source Gymnasium environment with WebSocket bridge to the
      real game engine, ensuring scoring fidelity.
  \end{itemize}
}

\plan{%
  Results preview (1 paragraph):
  \begin{itemize}
    \item Frontier LLMs beat a strong greedy baseline across all tested
      configs in metadata mode.
    \item Random agent defines a clear floor (score $\approx$ starting mana).
    \item Greedy agent provides a calibrated mid-point.
    \item Difficulty scales predictably with the number of active dimensions.
    \item \todo{Vision mode results once we have them.}
  \end{itemize}
}


% ======================================================================
\section{Related Work}
\label{sec:related}

Our work sits at the intersection of game-based AI evaluation, multi-objective
decision making, and vision-language model benchmarking. We organize prior work
into four areas.

\subsection{Game-Based Benchmarks for LLM and VLM Agents}

A growing body of work uses games to evaluate language and vision-language
models as interactive agents. SmartPlay~\citep{wu2024smartplay} was among the
first dedicated benchmarks, testing LLMs across six text-based games spanning
spatial reasoning, planning, and object dependencies. GameBench~\citep{costarelli2024gamebench}
extended this to nine game environments, each targeting a different strategic
reasoning axis, while GTBench~\citep{duan2024gtbench} introduced game-theoretic
structure with Monte Carlo Tree Search opponents across complete- and
incomplete-information settings. LLMArena~\citep{chen2024llmarena} and
GameArena~\citep{hu2025gamearena} further explored multi-agent competitive
evaluation and live human--AI interaction respectively.
DSGBench~\citep{tang2025dsgbench} shares our emphasis on multi-dimensional
evaluation, scoring agents across five specific capability dimensions within
complex strategy games.

On the vision side, lmgame-Bench~\citep{hu2026lmgamebench} wraps six
video/puzzle games in a unified Gymnasium API with perception scaffolds for
LLMs, while V-MAGE~\citep{zheng2025vmage} evaluates VLMs on arcade games
where all information is conveyed solely through screenshots.
Cradle~\citep{tan2025cradle} demonstrated general computer control via
screenshots in open-world games (Red Dead Redemption~2, Cities: Skylines),
and VideoGameBench~\citep{zhang2025videogamebench} tested VLMs on 23 classic
video games, finding that even the best model completes under 1\% of the
benchmark. GameWorld~\citep{ouyang2026gameworld} provides a browser-based
benchmark with state-verifiable scoring, comparing raw-input and
semantic-action interfaces.

Most directly relevant is BALROG~\citep{paglieri2025balrog}, which aggregates
six RL game environments (BabyAI, Crafter, TextWorld, MiniHack, NetHack) and
evaluates both LLMs and VLMs with text and image observations. Their central
finding---that several models perform \emph{worse} with visual
representations---motivates our controlled ``vision tax'' measurement. Our work
complements BALROG by focusing on a card game domain that tests
multi-dimensional strategic reasoning rather than grid-world navigation.

Earlier foundations include the Arcade Learning Environment~\citep{bellemare2013ale},
which established pixel-based game evaluation for RL; the NetHack Learning
Environment~\citep{kuttler2020nethack}, which supports multiple observation
modalities for the same underlying game; TextWorld~\citep{cote2018textworld}
for procedurally generated text adventures;
MineDojo~\citep{fan2022minedojo} and MineRL~\citep{guss2019minerl} for
Minecraft-based benchmarks; AgentBench~\citep{liu2024agentbench} for
multi-environment LLM agent evaluation; and
Atari-GPT~\citep{waytowich2024atarigpt}, which applied modern VLMs to classic
Atari games as zero-shot controllers.

\subsection{Card Game AI}

Card games pose distinct challenges for AI: imperfect information, stochastic
draws, and strategic depth within compact state spaces.
Pluribus~\citep{brown2019pluribus} achieved superhuman performance in six-player
no-limit Texas Hold'em using game-theoretic self-play, and
PokerBench~\citep{zhuang2025pokerbench} later evaluated LLMs on 11,000 poker
scenarios, finding that even GPT-4 achieves only 54\% accuracy.
The Hanabi Challenge~\citep{bard2020hanabi} proposed the cooperative card game
Hanabi as an AI benchmark emphasizing theory-of-mind reasoning under imperfect
information. More recently, \citet{wang2025llmcardgames} systematically
evaluated LLMs across eight diverse card games (DouDizhu, Mahjong, Uno, Gin
Rummy, Hold'em variants), finding that fine-tuning can approach strong baselines
but at the cost of general capability. RL approaches have also been applied to
complex collectible card games such as Hearthstone~\citep{hearthstone2023}.
Cicero~\citep{meta2022cicero} demonstrated human-level play in Diplomacy by
combining language models with strategic reasoning in a multi-player,
imperfect-information setting.

Yedi differs from these in two key respects: (1)~merge outcomes are determined
by the \emph{product} of per-dimension multipliers across 1--4 independent
cognitive dimensions, so one bad dimension ruins the compound gain; and (2)~the
benchmark supports both structured-metadata and screenshot observation modes for
the same underlying task, enabling controlled measurement of the perceptual
grounding cost.

\subsection{Multi-Objective Decision Making}

The multi-objective RL (MORL) literature addresses agents that must optimize
vector-valued rewards. \citet{hayes2022morl} provide a practical guide covering
utility functions, scalarization approaches, and Pareto-based methods for
sequential decision problems. Yedi's multiplicative scoring can be viewed as a
specific scalarization: $r = \prod_{d \in \text{active}} m_d$, where $m_d$ is
the per-dimension merge multiplier. Unlike additive scalarizations (where
strength in one dimension compensates for weakness in another), the
multiplicative form is dominated by the \emph{worst} factor---a property that
makes balanced competence across all dimensions essential.

This connects to resource-constrained planning in LLM agents.
AucArena~\citep{chen2024aucarena} tests LLMs on budget management and
adaptive strategy in auction settings, while \citet{huang2024planningsurvey}
survey LLM planning capabilities including task decomposition, plan selection,
and reflection---all of which Yedi requires within its dual-constraint
(mana budget + move budget) framework.

\subsection{Cognitive Benchmarks and Vision-Language Evaluation}

Standard LLM benchmarks---MMLU~\citep{hendrycks2021mmlu},
GPQA~\citep{rein2023gpqa}, MATH~\citep{hendrycks2021math},
BIG-bench~\citep{srivastava2023bigbench} and its hard
subset~\citep{suzgun2023bigbenchhard}---test knowledge and reasoning through
independent questions with additive scoring. A model scoring 90\% on each of
four subtasks achieves a 90\% average; a model achieving $0.9\times$ in each of
four Yedi dimensions scores $0.9^4 = 0.66\times$, dramatically penalizing any
weakness. Chollet's \emph{On the Measure of Intelligence}~\citep{chollet2019measure}
and the Abstraction and Reasoning Corpus (ARC) advocate for benchmarks that test
skill acquisition and generalization from cognitive-science priors. Yedi shares
this philosophy---each game configuration presents a novel combination of
dimension rules---but extends it to sequential decision-making under resource
constraints.

Recent work has begun applying psychometric frameworks to LLMs.
\citet{reddy2025chcllm} applied the Cattell--Horn--Carroll (CHC) theory of
cognitive abilities to frontier models, finding a ``catastrophic paradox''
where models excel at crystallized knowledge but fail psychometric scoring.
\citet{pellert2024aipsychometrics} demonstrated that standard psychometric
inventories can measure LLM psychological traits. Yedi can be viewed as a
\emph{behavioral} cognitive assessment that measures fluid reasoning (Gf),
quantitative reasoning (Gq), and visual processing (Gv) through observed
gameplay rather than self-report, potentially avoiding the validity concerns
these works raise.

For vision-language evaluation specifically, MMBench~\citep{liu2024mmbench},
MM-Vet~\citep{yu2024mmvet}, and SEED-Bench~\citep{li2024seedbench} test
fine-grained perceptual abilities in static settings.
SeeClick/ScreenSpot~\citep{cheng2024seeclick} and
ScreenQA~\citep{hsiao2025screenqa} focus on GUI grounding from screenshots.
OSWorld~\citep{xie2024osworld} and AndroidWorld~\citep{rawles2025androidworld}
test agents on real computer and mobile tasks, while
WebArena~\citep{zhou2024webarena} and
VisualWebArena~\citep{koh2024visualwebarena} provide web-based agent
benchmarks. VisualWebArena is particularly relevant as it compares agents
using accessibility trees (structured text) versus screenshots for the same
web tasks---directly analogous to our metadata-versus-vision design.
Our benchmark builds on the Gymnasium standard
interface~\citep{towers2024gymnasium} and, while Unity
ML-Agents~\citep{juliani2018unityml} provides RL integration for Unity games,
we use a WebSocket bridge to the WebGL build to enable browser-based deployment
and native screenshot capture.


% ======================================================================
\section{The Yedi Game}
\label{sec:game}

\plan{%
  Explain the game to a reader who has never seen it. This section should
  be self-contained. Use a figure showing the game board layout.
}

\subsection{Board and Actions}
\plan{%
  \begin{itemize}
    \item Board: 1 ``new'' slot, 5 build slots, sell action.
    \item Actions: DRAW (spawn card), MOVE/MERGE (place or combine),
      SELL (convert card mana to score). 37 discrete actions total.
    \item Two scarce resources: mana (spent on draws, recovered on sells)
      and moves (fixed budget per game, typically 100).
    \item Score = peak mana ever reached during the game.
  \end{itemize}
  \todo{Figure: game board screenshot with annotations.}
}

\subsection{Dimensions and Modes}
\plan{%
  \begin{itemize}
    \item Each card carries attributes across 4 dimensions: number, color,
      shape, word. Each dimension has multiple sub-modes that change the
      merge semantics.
    \item Number: add, multiply, GCD, sort, trigonometry, vector, interval.
    \item Color: additive mixing, subtractive mixing, grayscale, text-color.
    \item Shape: triangle rotation, rectangle, Kanizsa illusions, Hanoi,
      sphere, triple.
    \item Word: antonyms (verbs, nouns, adjectives), synonyms, grammar
      pairs, question-answer pairs, scrabble.
    \item 3 additional dimensions exist in the full game (Music/beat timing,
      Memory/hidden cards, Motor/drag accuracy) but are not the primary
      focus of this benchmark. Memory is partially supported (perfect\_memory
      ablation flag).
  \end{itemize}
  \todo{Table: all modes with brief description and neutral/identity card.}
}

\subsection{Merge Mechanics}
\label{sec:merge}
\plan{%
  This is the core mechanic and the key thing that makes Yedi interesting.
  Explain carefully.
  \begin{itemize}
    \item When two cards merge, each active dimension independently produces
      a multiplier applied to the target card's mana.
    \item Multiplier tiers: PERFECT (2.5$\times$), GREAT (2.0$\times$),
      GOOD (1.5$\times$), OK ($\sim$1.1$\times$), NEUTRAL (1.0$\times$),
      BAD (0.9$\times$), TERRIBLE (0.7$\times$).
    \item Multipliers are applied sequentially across dimensions:
      $\text{mana}_{\text{new}} = \text{mana}_{\text{old}} \times
      \prod_{d \in \text{active}} m_d$.
    \item Key insight: one BAD dimension ($m_d = 0.9$) washes out two
      GREAT dimensions ($2.0 \times 2.0 \times 0.9 = 3.6$ vs.\
      $2.0 \times 2.0 \times 2.0 = 8.0$). The worst dimension bounds
      the outcome.
    \item Each card has a 3-merge cap. After 3 merges the card must be sold.
    \item Worked example with concrete numbers showing the exponential
      growth from chained perfect merges vs.\ the collapse from one
      bad merge.
  \end{itemize}
  \todo{Figure: merge multiplier interaction diagram.}\\
  \todo{Table: concrete example of 3-merge chain across 1, 2, 3, 4 dims.}
}


% ======================================================================
\section{Benchmark Design}
\label{sec:benchmark}

\subsection{Environment}
\plan{%
  \begin{itemize}
    \item Gymnasium interface: \texttt{reset()}/\texttt{step(action)}
      returning \texttt{(obs, reward, terminated, truncated, info)}.
    \item Action space: \texttt{Discrete(37)} --- 0=draw, 1--30=move/merge,
      31--36=sell.
    \item Observation: metadata dict (mana, slots with card attributes,
      action mask) or screenshot (960$\times$540 PNG).
    \item WebSocket bridge to Unity WebGL ensures merge scoring uses the
      identical C\# code as the real game --- no re-implementation drift.
    \item Deterministic seeding for reproducibility.
  \end{itemize}
  \todo{Architecture diagram: Python env $\leftrightarrow$ WebSocket
    $\leftrightarrow$ Unity.}
}

\subsection{Configurations}
\plan{%
  \begin{itemize}
    \item 40+ predefined configurations spanning 1--4 active dimensions.
    \item Difficulty tiers: easy (1 dim, simple modes), medium (2 dims),
      hard (3--4 dims, complex modes).
    \item Each config specifies which dimensions are active and which
      sub-mode each uses.
    \item Standardized episode length (100 moves), starting mana (200),
      merge cap (3).
  \end{itemize}
  \todo{Table: representative configs used in experiments, with mode
    combination and difficulty tier.}
}

\subsection{Observation Modes}
\plan{%
  \begin{itemize}
    \item \textbf{Metadata mode}: structured text description of each slot's
      card attributes (Num=3, Color=Red, Shape=triangle-up, Word=help),
      mana values, merge counts. All information is explicit.
    \item \textbf{Vision mode}: 960$\times$540 screenshot of the game board.
      The model must extract card attributes visually. Same underlying game
      state, same action space.
    \item Running the same config in both modes isolates the ``vision tax''
      --- the performance drop attributable to perceptual grounding rather
      than strategic reasoning.
    \item \textbf{Merge preview ablation}: optionally expose the bridge's
      pre-computed merge scores to the agent. Measures whether the model's
      bottleneck is evaluation (computing merge quality) or planning
      (choosing which merge to execute given known quality).
  \end{itemize}
}

\subsection{Metrics}
\plan{%
  \begin{itemize}
    \item \textbf{Primary}: max mana (peak total mana during the episode).
      Single scalar, matches the human scoring system.
    \item \textbf{Secondary}: surplus (max mana $-$ starting mana), total
      reward, steps used, game-over rate.
    \item \textbf{Diagnostic}: merge tier distribution (how many PERFECT,
      GREAT, GOOD, BAD merges), sell timing (merges done when sold),
      wasted draws, preview adherence rate, mana trajectory over time.
    \item Report mean $\pm$ std over $N$ episodes per config.
  \end{itemize}
}


% ======================================================================
\section{Agents}
\label{sec:agents}

\subsection{Random Agent (Floor)}
\plan{%
  Uniformly samples from valid actions. Establishes the score floor.
  Expected score $\approx$ starting mana (200) since random merges are
  equally likely to help or hurt, while draws steadily drain mana.
}

\subsection{Greedy Agent (Calibrated Baseline)}
\plan{%
  \begin{itemize}
    \item Uses the bridge's pre-computed merge previews to pick the merge
      with the highest minimum gain across active dimensions.
    \item Decision priority: sell maxed cards $\to$ merge if min-gain
      $\geq$ threshold $\to$ place into empty slot $\to$ sell new card.
    \item No lookahead beyond one step. No inter-slot merges.
    \item Represents the performance ceiling of ``perfect local information,
      no planning.''
  \end{itemize}
}

\subsection{LLM Agent (Metadata Mode)}
\plan{%
  \begin{itemize}
    \item Receives structured text description of game state + valid actions.
    \item System prompt explains game rules, merge mechanics, scoring.
    \item Outputs a single action integer per turn.
    \item Tested with: \todo{list models --- Claude, GPT-4, open-source
      models at various scales.}
    \item Compact history trail provides last $N$ steps as context.
  \end{itemize}
}

\subsection{VLM Agent (Vision Mode)}
\plan{%
  \begin{itemize}
    \item Receives screenshot + minimal text (mana, step count, valid
      actions).
    \item Must extract card attributes (numbers, colors, shapes, words)
      from the rendered image.
    \item Same strategic reasoning required as metadata mode, plus visual
      grounding.
    \item \todo{Tested with: Claude vision, GPT-4V, open-source VLMs.}
  \end{itemize}
}


% ======================================================================
\section{Experiments}
\label{sec:experiments}

\plan{%
  This section will be filled in as we run more benchmarks.
  Structure the results around 3 research questions.
}

\subsection{RQ1: Do LLMs exhibit multi-dimensional strategic reasoning?}
\plan{%
  \begin{itemize}
    \item Compare LLM agents vs.\ greedy vs.\ random across single-dim
      and multi-dim configs.
    \item Hypothesis: LLMs should beat greedy on multi-dim configs (where
      planning and trade-off reasoning matter) even if greedy has perfect
      local merge scores.
    \item Metric: max mana, broken down by number of active dimensions.
    \item \todo{Run matrix: 3+ LLMs $\times$ 8+ configs $\times$ 5 episodes.}
  \end{itemize}
  \todo{Table: agent $\times$ config score matrix.}\\
  \todo{Figure: bar chart grouped by dimension count.}
}

\subsection{RQ2: What is the vision tax?}
\plan{%
  \begin{itemize}
    \item Run the same LLM in metadata mode and vision mode on identical
      configs.
    \item The score difference isolates perceptual grounding cost from
      strategic reasoning.
    \item Break down by dimension type: color (requires seeing actual RGB),
      shape (spatial), number (OCR), word (text reading).
    \item Hypothesis: vision tax is smallest for number/word (essentially
      OCR) and largest for color/shape (genuine visual understanding).
    \item \todo{Run matrix: 2+ VLMs $\times$ same configs $\times$
      metadata and vision $\times$ 5 episodes.}
  \end{itemize}
  \todo{Table: metadata score vs.\ vision score, with delta column.}\\
  \todo{Figure: paired comparison per config.}
}

\subsection{RQ3: How does performance scale with model capability?}
\plan{%
  \begin{itemize}
    \item Test models across a capability range: small open-source
      ($\sim$7B), medium ($\sim$70B), large ($\sim$120B+), frontier
      (Claude, GPT-4).
    \item Does Yedi score correlate with general benchmarks (MMLU, HumanEval)?
      Or does it measure something orthogonal?
    \item Hypothesis: multi-dim configs should differentiate models more
      than single-dim configs (harder task = more headroom).
    \item \todo{Run matrix: 5+ models $\times$ easy/medium/hard configs
      $\times$ 5 episodes.}
  \end{itemize}
  \todo{Figure: scaling curve (model capability vs.\ Yedi score).}\\
  \todo{Scatter: Yedi score vs.\ MMLU / HumanEval.}
}

\subsection{Ablation: Merge Previews}
\plan{%
  \begin{itemize}
    \item With previews: agent receives pre-computed merge quality scores.
      Tests pure planning (``given that you know the quality of every
      possible merge, can you pick the right sequence?'').
    \item Without previews: agent must compute merge quality from card
      attributes and game rules. Tests evaluation + planning.
    \item Delta measures the ``evaluation bottleneck'' --- how much
      performance is lost because the model can't accurately score merges.
    \item \todo{Run same model with and without previews on same configs.}
  \end{itemize}
}


% ======================================================================
\section{Analysis}
\label{sec:analysis}

\plan{%
  Deeper dives beyond the headline numbers.
}

\subsection{Merge Quality Distribution}
\plan{%
  \begin{itemize}
    \item For each agent, plot the distribution of merge tiers chosen
      (PERFECT, GREAT, GOOD, BAD).
    \item Do LLMs make fewer BAD merges than greedy? (Greedy has perfect
      local info but no planning; LLMs might avoid traps greedy walks into.)
    \item Mana trajectory plots: how does mana evolve over the 100-step
      episode for each agent type?
  \end{itemize}
  \todo{Figure: merge tier histograms per agent.}\\
  \todo{Figure: mana trajectory curves (mean $\pm$ std).}
}

\subsection{Failure Modes}
\plan{%
  \begin{itemize}
    \item Qualitative analysis of LLM mistakes from trace logs.
    \item Common failure patterns: single-slot fixation, forgetting to sell,
      bad-merge spam, ignoring one dimension.
    \item Do different models exhibit different failure modes?
    \item \todo{Pull concrete examples from trace.jsonl files.}
  \end{itemize}
}

\subsection{Dimension-Specific Performance}
\plan{%
  \begin{itemize}
    \item Per-dimension merge quality breakdown: is the model better at
      number reasoning than color reasoning?
    \item Radar chart: one axis per dimension, comparing agents.
    \item Which dimension is the bottleneck for each model?
  \end{itemize}
  \todo{Figure: radar chart comparing agents across dimensions.}
}


% ======================================================================
\section{Limitations}
\label{sec:limitations}

\plan{%
  Be upfront. Reviewers will find these anyway.
  \begin{itemize}
    \item \textbf{Visual complexity}: the game UI is clean and structured.
      Vision mode tests ``can you read a game board'' more than ``can you
      understand a complex visual scene.'' The vision tax may underestimate
      VLM limitations on messier real-world visual reasoning.
    \item \textbf{Game-specific knowledge}: models that have encountered
      Yedi-like rules in training data may have an advantage. We mitigate
      this with the diversity of mode combinations and the novelty of the
      specific merge mechanics, but cannot fully rule it out.
    \item \textbf{Bridge fragility}: the WebSocket bridge to Unity WebGL
      introduces engineering complexity and occasional disconnections.
      Episode-level error recovery is implemented but adds noise.
    \item \textbf{Evaluation cost}: each LLM episode requires $\sim$100
      API calls. Large-scale evaluation across many models and configs is
      expensive. We report results with limited episode counts and
      acknowledge the statistical uncertainty.
    \item \textbf{Single-player}: Yedi is not adversarial. The benchmark
      tests optimization against a fixed environment, not adaptation to
      an opponent. Strategic depth is bounded by the merge mechanics
      rather than an open-ended opponent.
    \item \textbf{Action space simplicity}: 37 discrete actions is not a
      complex action space. The challenge is in state evaluation and
      planning, not in action selection or motor control.
  \end{itemize}
}


% ======================================================================
\section{Conclusion}
\label{sec:conclusion}

\plan{%
  3 paragraphs:
  \begin{enumerate}
    \item We presented Yedi, a benchmark for multi-dimensional strategic
      reasoning. The multiplicative interaction across dimensions means
      models must optimize the worst-case factor, not the average ---
      a capability not tested by existing single-axis game benchmarks.
    \item Key findings: \todo{summarize RQ1--RQ3 results in 2--3 sentences
      each once experiments are complete.}
    \item Future work: RL training on Yedi, the 3 remaining dimensions
      (Music, Memory, Motor), multi-player/adversarial variants,
      curriculum learning across the difficulty ladder.
  \end{enumerate}
}


% ======================================================================
% Optional: Appendix
% ======================================================================
\appendix

\section{Full Configuration Table}
\plan{%
  Complete list of all 40+ benchmark configurations with active dimensions,
  sub-modes, and difficulty tier.
}

\section{Prompt Templates}
\plan{%
  The system prompt and per-turn user message templates used for LLM/VLM
  agents. Included for reproducibility.
}

\section{Per-Dimension Merge Rules}
\plan{%
  Detailed scoring tables for each dimension and sub-mode. Reference for
  readers who want to understand exactly how each multiplier tier is
  determined.
}


% ======================================================================
\bibliographystyle{plainnat}
\bibliography{references}
\plan{%
  \todo{Create references.bib with citations for:
    GamingAgent, V-MAGE, Cradle, WebArena, VisualWebArena,
    Unity ML-Agents, Gymnasium/Farama, ARC, BIG-bench,
    MORL survey, Fitts' Law, CHC cognitive theory,
    Claude/GPT-4 technical reports, MineRL, NetHack LE.}
}

\end{document}
